\chapter{Principais trechos do firmware do bracelete}
\label{ap:codigo}

\lstset{
  basicstyle=\ttfamily\footnotesize,
  columns=fullflexible,
  keepspaces=true,
  frame=single,
  breaklines=true,
  showstringspaces=false,
  language=C++
}

Os trechos a seguir foram extraídos do firmware do bracelete (\texttt{bracelete\_ble.ino}). O código completo, o firmware da baliza e a interface web estão no repositório do projeto.

\section{Recebimento dos anúncios e suavização do RSSI}

A cada anúncio recebido, o bracelete verifica se o nome começa com \texttt{BALIZA-} e atualiza o RSSI suavizado da baliza correspondente (Equação~\ref{eq:ema}).

\begin{lstlisting}
void onResult(BLEAdvertisedDevice dev) override {
  if (!dev.haveName()) return;
  String full = dev.getName();
  if (!full.startsWith(BEACON_PREFIX)) return;
  String name = full.substring(strlen(BEACON_PREFIX));
  int rssi = dev.getRSSI();
  ...
  if (slot < 0) {                       // baliza nova
    strlcpy(beacons[slot].name, name.c_str(), sizeof(beacons[slot].name));
    beacons[slot].rssi = rssi;
  } else {                              // media movel exponencial
    beacons[slot].rssi += RSSI_ALPHA * (rssi - beacons[slot].rssi);
  }
  beacons[slot].lastSeen = now;
}
\end{lstlisting}

\section{Estimativa de distância e histerese}

\begin{lstlisting}
int rssiToCm(float rssi) {              // Equacao (1), em centimetros
  return 100 * pow(10, (RSSI_1M - rssi) / (10 * PATH_LOSS_N));
}

const int LIMITS[] = {400, 200, 100};   // faixas: 4 m, 2 m e 1 m

int levelWithHysteresis(int cm, int current) {
  int level = 0;
  for (int i = 0; i < 3; i++) if (cm <= LIMITS[i]) level = i + 1;
  if (level < current && cm <= LIMITS[current - 1] * 1.25) return current;
  return level;
}
\end{lstlisting}

\section{Característica sem escrita e envio ao celular}

A característica é criada apenas com leitura e notificação, o que garante o fluxo em sentido único. No laço principal, a lista de balizas é montada da mais próxima para a mais distante e enviada a cada 500~ms.

\begin{lstlisting}
distChar = service->createCharacteristic(
    DIST_CHAR_UUID,
    BLECharacteristic::PROPERTY_READ | BLECharacteristic::PROPERTY_NOTIFY);

// laco principal: "<demo>|Nome:cm;Nome:cm;..."
String payload = String(demoMode ? 1 : 0) + "|";
if (connected && now - lastNotify >= NOTIFY_MS) {
  size_t maxLen = server->getPeerMTU(server->getConnId()) - 3;
  for (int i = 0; i < n; i++) {
    String item = String(i ? ";" : "") + points[i].name + ":" + points[i].cm;
    if (payload.length() + item.length() > maxLen) break;
    payload += item;
  }
  distChar->setValue(payload.c_str());
  distChar->notify();
  lastNotify = now;
}
\end{lstlisting}
